% =====================================================================
%  Annotated REVTeX 4.2 skeleton for a Physical Review manuscript.
%
%  Swap the class options for the target journal:
%    PRL      : [aps,prl,preprint,superscriptaddress]
%    PRB      : [aps,prb,preprint,superscriptaddress]
%    PRX      : [aps,prx,preprint,superscriptaddress]
%    APL (AIP): [aip,apl,preprint]        -- superscript citations
%    JAP (AIP): [aip,jap,preprint]
%
%  Submit with `preprint`. Compile locally with `reprint` to check the
%  true two-column page count before you claim it fits.
% =====================================================================

\documentclass[
  aps,prl,
  preprint,              % <- switch to `reprint` to check page count
  amsmath,amssymb,
  superscriptaddress,
  showpacs,
  nofootinbib
]{revtex4-2}

\usepackage{graphicx}
\usepackage{bm}          % \bm{B} for vectors -- APS style is bold, not arrows
\usepackage{siunitx}     % \qty{5}{\kelvin}, \qty{42(5)}{\watt\per\meter\per\kelvin}
\usepackage{dcolumn}     % decimal-aligned table columns
\usepackage{booktabs}    % \toprule \midrule \bottomrule -- no vertical rules
\usepackage[colorlinks=true,allcolors=blue]{hyperref}

\sisetup{separate-uncertainty = true, range-units = single}

% --- Keep macros few and mnemonic; APS asks that they be minimal ------
\newcommand{\kperp}{k_{\perp}}
\newcommand{\Tc}{T_{\mathrm{c}}}

\begin{document}

% ---------------------------------------------------------------------
%  TITLE: declarative, specific, no new terminology, no unfamiliar
%  acronyms. Prefer "Observation of X in Y" over "A study of X".
% ---------------------------------------------------------------------
\title{Declarative statement of the principal result}

\author{A.~Author}
\email{corresponding@institution.edu}
\affiliation{Department of Physics, University, City 12345, Country}

\author{B.~Author}
\affiliation{Second Institution, City 67890, Country}

\date{\today}

% ---------------------------------------------------------------------
%  ABSTRACT: one unlabeled paragraph. PRL: 600 characters INCLUDING
%  SPACES (~90-100 words). Five beats:
%    (1) context  (2) gap  (3) what we did
%    (4) THE RESULT WITH A NUMBER AND UNCERTAINTY  (5) consequence
%  No citations, no undefined acronyms, no figure references.
% ---------------------------------------------------------------------
\begin{abstract}
  ...
\end{abstract}

\maketitle   % <- NOTE: comes AFTER the abstract in REVTeX

% ---------------------------------------------------------------------
%  A PRL runs as continuous prose; sections below are for an Article.
%  Delete the \section commands for a Letter.
% ---------------------------------------------------------------------

\section{Introduction}
% Move 1: establish the phenomenon and why it matters (1-2 sentences)
% Move 2: what is established, WITH NUMBERS and selective citations
% Move 3: the gap -- no measurement / unexplored regime / theory-experiment
%         discrepancy / conflicting results / method limitation
% Move 4: "Here we ..." + the headline result with its number
% Move 5: roadmap (Articles only -- omit in a Letter)

\section{Theory and model}
% Assumptions and validity range first, then the governing equations.
% Define every symbol. Recover known limits. State conventions once:
% units, metric signature, Fourier sign, natural units.

Substituting the interaction term into the Hamiltonian gives
\begin{align}
  \mathcal{H} &= \sum_i \epsilon_i c_i^{\dagger} c_i
                 + \sum_{ij} t_{ij}\, c_i^{\dagger} c_j , \label{eq:ham}
\end{align}
where $\epsilon_i$ is the on-site energy, $t_{ij}$ the hopping amplitude,
and $c_i^{\dagger}$ ($c_i$) creates (annihilates) an electron on site $i$.
% ^ Note the comma after the displayed equation: it is part of this sentence.
% Reference it later as Eq.~(\ref{eq:ham}) or \eqref{eq:ham} -- never "(1)".

\section{Experimental methods}
% Move 1: general overview before any parameter
% Move 2: specific parameters (quantities, temperatures, durations, geometry)
% Move 3: justify the choices
% Move 4: indicate care taken (calibration, controls, repeats, convergence)
% Move 5: acknowledge method limitations WITH BOUNDS
% Past tense throughout for what was done.

\section{Results}
% One prose block per figure:
%   invitation -> salient feature -> quantify with uncertainty -> compare
% Present tense for what the figure shows; past for what you found.

\begin{figure}
  \includegraphics[width=\columnwidth]{fig1.pdf}
  % Design at 8.6 cm (single column). Fonts 7-9 pt AT THAT SIZE.
  \caption{\label{fig:one}%
    Title phrase stating what the figure shows. (a) What is plotted, with
    symbols and conditions. (b) Second panel. Solid line: model; markers:
    data. Error bars denote one standard deviation over five thermal
    cycles; where not visible they are smaller than the markers.}
\end{figure}

\begin{table}
  \caption{\label{tab:budget}%
    Systematic uncertainty budget. Caption goes ABOVE the table.}
  \begin{ruledtabular}
  \begin{tabular}{lD{.}{.}{2}}
    Source & \multicolumn{1}{c}{Contribution (\%)} \\
    \colrule
    Thermometer calibration      & 4.0 \\
    Radiative loss correction    & 2.5 \\
    Sample geometry              & 3.1 \\
    \colrule
    Total (quadrature)           & 5.8 \\
  \end{tabular}
  \end{ruledtabular}
\end{table}

\section{Discussion}
% Restate the finding -> mechanism -> quantitative comparison with prior
% work -> rule out the two leading alternatives -> limitations with bounds
% -> implications. Calibrate every verb to the evidence.

\section{Conclusions}
% One paragraph. Claim + number, significance, next question.
% No new results, no new references.

\begin{acknowledgments}
  Funding, facilities, and individuals. (Anonymize this block if the
  journal uses anonymous review.)
\end{acknowledgments}

% --- Data and code availability --------------------------------------
% Name the repository and give an archived DOI; "available on reasonable
% request" alone is increasingly rejected.

\appendix

\section{Derivation of Eq.~(\ref{eq:ham})}
% For a PRL, material that specialists need goes in End Matter (a specific
% APS feature placed after the references) rather than a standard appendix.
% Check the current APS markup for End Matter before using it.

\bibliography{refs}   % use the publisher's .bst; never hand-format entries

\end{document}
